% TOP_PROBLEM: {"problemId": "problem.hilbert-sixteenth-problem-second-part", "canonicalTitle": "Hilbert's Sixteenth Problem, Second Part", "releaseRank": 40, "primaryDomain": "probability_ergodic_dynamics", "primaryDomainLabel": "Probability, ergodic theory & dynamics", "displayStatus": "open", "releaseStatus": "open"}
% !TeX program = lualatex
% Definition and sources only; this file is not an LLM solution attempt.
\documentclass[11pt]{article}
\usepackage[margin=25mm]{geometry}
\usepackage{fontspec}
\setmainfont{Latin Modern Roman}
\usepackage{amsmath,amssymb,mathtools}
\usepackage{unicode-math}
\setmathfont{Latin Modern Math}
\usepackage{xurl,hyperref}
\hypersetup{colorlinks=true,urlcolor=blue}
\setcounter{secnumdepth}{0}
\setlength{\parindent}{0pt}
\setlength{\parskip}{5pt}
\setlength{\emergencystretch}{3em}
\begin{document}
\section*{\texorpdfstring{Hilbert's Sixteenth Problem, Second Part}{Hilbert's Sixteenth Problem, Second Part}}\label{40}
\subsection{Definitions and mathematical statement}
A planar polynomial vector field of degree at most $n$ is
\[
 \dot x=P(x,y),\qquad\dot y=Q(x,y),\qquad P,Q\in{\mathbb{R}}[x,y],\quad\max(\deg P,\deg Q)\le n.
\]
A limit cycle is an isolated periodic orbit, not an arbitrary member of a continuous family of periodic trajectories. Let $H(n)$ be the supremum of the number of limit cycles over this class. The problem asks for finiteness of $H(n)$, its maximal value, and the possible relative configurations of these cycles for each degree. A bound depending on the individual coefficients is not a uniform degree bound.
\par\smallskip\textit{Formulation and concept references:} \href{https://recercat.cat/bitstream/handle/2072/532013/MelnikovMethod.pdf?sequence=1}{[S1]}.

\subsection{Short English statement}
How many isolated periodic motions can a planar polynomial differential equation have, and how can those cycles be arranged?
\subsection{Sources}
\begin{itemize}
\item[S1]Luiz Fernando da Silva Gouveia and Joan Torregrosa, The local cyclicity problem. Melnikov method using Lyapunov constants, Proceedings of the Edinburgh Mathematical Society 65 (2022), 356-375, Introduction, p. 356.
\url{https://recercat.cat/bitstream/handle/2072/532013/MelnikovMethod.pdf?sequence=1}.
\textit{Formulation source.}
\item[S2]A note on a recent attempt to solve the second part of Hilbert's 16th Problem.
\url{https://arxiv.org/abs/2411.09594}.
\textit{Further reference; not a proof endorsement.}
\item[S3]Pesquisadores da Unesp propõem solução para desafio matemático em aberto há mais de um século.
\url{https://jornal.unesp.br/2024/09/24/pesquisadores-da-unesp-propoem-solucao-para-desafio-matematico-em-aberto-ha-mais-de-um-seculo/}.
\textit{Further reference; not a proof endorsement.}
\end{itemize}
\end{document}
